%=====================================================================
\section{Introduction}
\label{sec:intro}
%=====================================================================

As language models (LMs) are deployed in more applications, the ability to erase specific knowledge from them becomes important for safety and compliance \citep{eldan2023harry,li2024wmdp}.
A prominent approach is \emph{persistent} concept erasure, which removes knowledge of a target concept by modifying the model's weights \citep{gurarieh2025precise,ashuach2026crisp}.
Recently, \citet{suslik2026ember} showed that such methods overlook a part of the model that stores knowledge: the token embeddings.
Their method, \ember{}, decomposes the embeddings of concept-related tokens into sparse features, selects the features of the target concept with an LLM judge, and subtracts them from a small set of tokens.\footnote{Our code is available at \url{https://github.com/guretskiysemyon/nlp_proj}.}
Combined with MLP-based methods, it substantially improves robustness to relearning.

\begin{figure}[t]
\centering
\begin{tikzpicture}[
  box/.style={draw, rounded corners=2pt, align=center, font=\scriptsize, text width=2.85cm,
              minimum height=0.86cm, inner sep=2pt},
  eb/.style={box, draw=embergrey, fill=embergrey!7},
  cb/.style={box, draw=cinderorange, fill=cinderorange!7, line width=0.6pt},
  sl/.style={font=\scriptsize\bfseries, align=right, text width=1.15cm, anchor=east},
  hd/.style={font=\footnotesize\bfseries},
  arr/.style={-{Stealth[length=3pt]}, semithick, draw=black!45}]
  \node[eb] (e1) at (0,0) {Sparse MF of raw embeddings;\\ 1\% of tokens per feature};
  \node[eb, below=0.22cm of e1] (e2) {$\rho > \tau$; LLM judge reads\\ a one-line summary};
  \node[eb, below=0.22cm of e2] (e3) {$\e_t - \delta\,\cvec_t$, with $\delta$\\ up to 200};
  \node[eb, below=0.22cm of e3] (e4) {efficacy, MMLU,\\ Alpaca coherence};
  \node[eb, below=0.22cm of e4] (e5) {English only; one fit\\ per language otherwise};
  \node[cb, right=0.25cm of e1] (c1) {\textbf{sharpen} toward the concept;\\ \textbf{one global} sparsity budget};
  \node[cb, right=0.25cm of e2] (c2) {judge \textbf{or} lexicon match\\ (one LLM call per concept)};
  \node[cb, right=0.25cm of e3] (c3) {same tokens; concept part\\ \textbf{relative to ordinary text}};
  \node[cb, right=0.25cm of e4] (c4) {+ \textbf{collateral}, edit size,\\ \textbf{word integrity}};
  \node[cb, right=0.25cm of e5] (c5) {\textbf{one fit for eight}\\ \textbf{languages}};
  \node[sl, left=0.12cm of e1] {Decompose};
  \node[sl, left=0.12cm of e2] {Identify};
  \node[sl, left=0.12cm of e3] {Erase};
  \node[sl, left=0.12cm of e4] {Evaluate};
  \node[sl, left=0.12cm of e5] {Languages};
  \node[hd, text=embergrey, above=0.05cm of e1] {\ember{}};
  \node[hd, text=cinderorange, above=0.05cm of c1] {\cinder{} (ours)};
  \foreach \a/\b in {e1/e2,e2/e3,c1/c2,c2/c3} \draw[arr] (\a) -- (\b);
  \draw[arr, dashed] (e3) -- (e4);
  \draw[arr, dashed] (c3) -- (c4);
\end{tikzpicture}
\caption{\ember{} erases a concept by decomposing token embeddings into sparse features, identifying the concept's features, and subtracting them from concept tokens. \cinder{} keeps this framework and changes how the concept is found, verified and edited, so that quiet concepts become the strongest direction of the matrix. One \cinder{} fit serves eight languages, where \ember{} needs a fit per language. We also extend the evaluation so that destroying the edited words no longer counts as erasing the concept.}
\label{fig:overview}
\end{figure}

However, an embedding edit can only remove what its features capture, and how reliably \ember{} finds a concept's features has not been studied.
Its authors note that embeddings may not be an effective target for some concepts, that the tokens selected for editing may include semantically unrelated ones, and that the edit only applies to English tokens \citep[Limitations]{suslik2026ember}.

In this work, we hypothesize that these limitations stem in part from the geometry in which \ember{} searches for features.
Specifically, sparse matrix factorization minimizes reconstruction error, so it spends its features on directions that many tokens share, such as magic, writing or books, while a concept such as Harry Potter is carried by a few rare tokens and explains little of the error.
We call such a concept \emph{quiet}: in \ember{}'s matrix, eleven generic features outrank the Harry Potter feature on \ember{}'s own concept score, and the feature's tail holds tokens such as Nestl\'e, Lego and Trump (Figure~\ref{fig:hp}).
Two more design choices compound the problem.
Every feature keeps the same 1\% of the tokens regardless of how broad its concept is, and the judge accepts a feature based on a one-sentence summary that often abstracts the concept away.
When no feature is accepted, nothing is erased.

To tackle this gap, we introduce \cinder{} (Contrastive INformed Directional Embedding Removal; Figure~\ref{fig:overview}), which keeps \ember{}'s localize-then-edit framework and changes how the concept is found.
\cinder{} whitens the embeddings against ordinary text and sharpens them toward the concept's page, so that what distinguishes the concept becomes the strongest direction in the matrix.
It shares one sparsity budget across all features, so that each feature takes as many tokens as the data supports, and it verifies features both bottom-up, with \ember{}'s judge, and top-down, against a concept lexicon generated once per concept.
Because the sharpened space is centred on ordinary text, \cinder{}'s edit removes only what distinguishes the concept; an uncentred variant, \cinderz{}, also removes part of the component that all tokens share, as \ember{}'s edit does.

Sharpening also makes multilingual erasure practical, which \ember{} leaves to future work \citep[Limitations]{suslik2026ember}.
Because \ember{} edits the rows of English tokens, a concept it erases in English largely survives in other languages: on Gemma, edited models retain \expnum{0.24} more of the concept's knowledge in German, Spanish, French and Russian than in English.
Running \ember{} separately in every language multiplies its cost by the number of languages, and fitting it once on multilingual text mostly fails, yielding a concept feature for only \expnum{10 of 38} concepts.
We expect the concept to be quieter still in a multilingual matrix, where language identity adds loud directions of its own.
\cinder{} whitens against a background with equal text in every language, which removes these directions, and sharpens toward the concept's pages in all languages.
A single \cinder{} fit over eight languages then erases as well as eight separate \ember{} runs, with one factorization and one judging pass instead of eight; to our knowledge, this is the first multilingual erasure at the embedding level.

We assess \cinder{} through experiments on Gemma-2-2B-it \citep{gemma2024gemma2} and, across eight languages, Llama-3.1-8B-Instruct \citep{grattafiori2024llama3}.
Notably, evaluating it exposed a blind spot in the standard protocol.
\ember{}'s tuning score prefers edits up to a hundred times longer than the embeddings they change, which destroy the edited words rather than remove the concept from them.
Because the concept questions contain the edited words, a destroyed word counts as erased knowledge, while MMLU and Alpaca coherence rarely encounter it.
We therefore extend the evaluation with \emph{collateral} questions that use the edited words in their other meanings (e.g., Prince Harry when erasing Harry Potter), one scale of edit size shared by all methods, and tests of whether the edited words survive.

Our results show that once collateral damage is charged for, every method moves from token-destroying edits to edits about as long as the token, and that embedding erasure blocks access to a concept through its words rather than removing it from the weights.
Within this regime, \cinder{} erases concepts that \ember{} misses (\expnum{14} vs.\ \expnum{3} runs in which only one of the methods edits), and its direction removes more at equal edit size (\expnum{$+0.15$} on the combined score, \expnum{10 of 11} held-out concepts), while at each method's own operating point the two tie.
Across eight languages, a single \cinder{} fit matches running \ember{} separately in each language (\expnum{0.23} vs.\ \expnum{0.25}) at one eighth of the selection cost, and finds concept features for twice as many concepts as a single multilingual \ember{} fit (\expnum{21} vs.\ \expnum{10} of 38).

Finally, we analyze how sharpening changes the embedding space.
We find that the model's first layer reads one direction, the mean of all embeddings, far more strongly than others, and that this shared direction, not the concept, explains why \ember{}'s edits are read more strongly than \cinder{}'s.
\pending{one or two sentences from the embedding-space analysis (\S\ref{sec:analysis}): the concept direction's energy rank, near-orthogonal features with meaningful angles, widths that follow the data, and the circularity control.}

To conclude, our work (1) diagnoses why embedding-level erasure misses quiet concepts, (2) introduces \cinder{}, which finds and erases them more reliably, (3) extends embedding-level erasure to many languages with a single fit, at one eighth of the selection cost of erasing each language separately, and (4) shows that the standard evaluation of embedding erasure rewards destroying tokens, and proposes a collateral-aware alternative that needs no LLM judge inside the tuning loop.
Overall, our findings suggest that where a concept is looked for matters as much as how it is removed, and that an embedding edit must be judged by what it does to the words it edits.

